\documentclass[11pt,a4paper]{article}

\usepackage[T1]{fontenc}
\usepackage[utf8]{inputenc}
\usepackage[italian]{babel}
\usepackage{lmodern}
\usepackage{amsmath}
\usepackage{microtype}
\usepackage[margin=2.5cm]{geometry}
\usepackage{xcolor}
\usepackage{booktabs}
\usepackage{tabularx}
\usepackage{caption}
\usepackage{listings}
\usepackage{tikz}
\usetikzlibrary{positioning,arrows.meta,shapes.misc,fit,calc}
\usepackage[hidelinks]{hyperref}
\setlength{\emergencystretch}{3em}

% ---------------------------------------------------------------------------
% Macro di servizio
% ---------------------------------------------------------------------------
\newcommand{\todo}[1]{\textcolor{red}{\textbf{[TODO: #1]}}}
\newcommand{\dq}{\textquotedbl}            % virgolette dritte nelle query
\newcommand{\code}[1]{\texttt{#1}}
\newcommand{\field}[1]{\texttt{#1}}

\definecolor{accent}{HTML}{1F6F8B}
\definecolor{light}{HTML}{E8F1F5}
\definecolor{warm}{HTML}{C8553D}
\definecolor{codebg}{HTML}{F6F8FA}

\lstset{
  language=Java,
  basicstyle=\ttfamily\small,
  keywordstyle=\color{accent}\bfseries,
  commentstyle=\color{gray}\itshape,
  backgroundcolor=\color{codebg},
  frame=single, rulecolor=\color{gray!40},
  columns=fullflexible, keepspaces=true,
  showstringspaces=false, breaklines=true,
  xleftmargin=4pt, xrightmargin=4pt
}

\tikzset{
  box/.style={draw=accent, fill=light, rounded corners=3pt, align=center,
              minimum height=8mm, inner sep=4pt, font=\small},
  store/.style={draw=warm, fill=warm!10, rounded corners=3pt, align=center,
                minimum height=8mm, inner sep=4pt, font=\small},
  tok/.style={draw=accent, fill=white, inner sep=3pt, font=\ttfamily\small},
  arr/.style={-{Stealth[length=2.5mm]}, thick, accent},
}

\title{Homework 3: indicizzazione e ricerca di file di testo con Lucene}
\author{Carmine Adduce \\ \small Matricola 589659 \\
            \small E-mail car.adduce@stud.uniroma3.it \\
            \small Corso di Ingegneria dei Dati}
\date{Novembre 2026}

\begin{document}
\maketitle

\begin{abstract}
Questa relazione descrive un programma Java che indicizza con Apache Lucene i file
\code{.txt} di una cartella, usando due campi (nome del file e contenuto) con un
analyzer dedicato ciascuno, e che permette di interrogare l'indice da console con
una sintassi semplice. Il sistema è stato sviluppato e provato con 21 query su un
corpus bilingue (italiano e inglese) di 202 file, poi misurato su un corpus più
grande di 2\,810 file (circa 34~MB) costruito con uno script riproducibile da
Wikipedia: l'indicizzazione richiede circa 2,3 secondi e il tempo cresce quasi
linearmente con il numero di file. La qualità è valutata con precision e recall
rispetto a una verità lessicale indipendente da Lucene (e controllata con
\code{grep}).
\end{abstract}

\noindent\textbf{Repository:}
\url{https://github.com/adducec03/lucene-file-indexer}

% ===========================================================================
\section{Obiettivo e tecnologie}
% ===========================================================================

La consegna chiedeva di:
\begin{enumerate}
  \item indicizzare i file \code{.txt} di una cartella su due campi, nome e
        contenuto, scegliendo per ciascuno un analyzer appropriato;
  \item leggere da console una query che inizia con la parola chiave del campo
        (\field{nome} o \field{contenuto}), seguita da termini eventualmente tra
        virgolette per esprimere una \emph{phrase query}, e stamparne i risultati;
  \item provare il sistema con una decina di query diverse.
\end{enumerate}

\noindent Il progetto usa Java~21, Maven (con il wrapper \code{./mvnw}, che non
richiede installazioni) e \textbf{Lucene~10.3.1}. Le dipendenze sono \code{lucene-core}, \code{lucene-analysis-common} e
\code{lucene-queryparser}.

\noindent Oltre alla consegna, il sistema è stato provato anche su un corpus molto più
grande di quello usato per lo sviluppo, per misurare come crescono i tempi di
indicizzazione e per valutare la qualità dei risultati con metriche (precision e
recall), non solo contando i file trovati.

% ===========================================================================
\section{Architettura}
% ===========================================================================

Il codice è nel package \code{it.homework.lucene} ed è diviso in cinque classi,
ciascuna con un solo compito (tabella~\ref{tab:classi}), più una sesta classe,
\code{Evaluator}, che serve solo per la valutazione (sezione~\ref{sec:valutazione}). I due flussi principali,
indicizzazione e ricerca, condividono lo stesso analyzer: è essenziale, perché i
termini della query devono essere trattati esattamente come quelli dell'indice
(figura~\ref{fig:architettura}).

\begin{table}[!htbp]
\centering
\begin{tabularx}{\linewidth}{@{}lX@{}}
\toprule
\textbf{Classe} & \textbf{Responsabilità} \\
\midrule
\code{Main} & comandi \code{index} e \code{search}, ciclo interattivo su console \\
\code{Analyzers} & nomi dei campi e analyzer per campo (\code{PerFieldAnalyzerWrapper}) \\
\code{FileIndexer} & visita la cartella, crea un \code{Document} per file, misura i tempi \\
\code{QueryBuilder} & traduce la riga digitata in una \code{Query} di Lucene \\
\code{FileSearcher} & apre l'indice in lettura, esegue la query e stampa i risultati \\
\code{Evaluator} & (solo valutazione) calcola precision e recall confrontando i risultati con una verità lessicale \\
\bottomrule
\end{tabularx}
\caption{Le classi del progetto.}
\label{tab:classi}
\end{table}

\begin{figure}[!htbp]
\centering
\begin{tikzpicture}
  % --- indicizzazione (riga in alto)
  \node[store] (docs)  at (0,0)    {cartella\\\code{data/docs}};
  \node[box]   (idx)   at (3.7,0)  {\code{FileIndexer}};
  \node[box]   (iw)    at (7.5,0)  {\code{IndexWriter}};
  \node[store] (index) at (11,0)   {indice\\Lucene};
  \draw[arr] (docs) -- node[above, font=\scriptsize]{file .txt} (idx);
  \draw[arr] (idx) -- node[above, font=\scriptsize]{Document} (iw);
  \draw[arr] (iw) -- node[above, font=\scriptsize]{commit} (index);

  % --- analyzer condiviso (al centro)
  \node[box, fill=warm!15, draw=warm] (an) at (5.6,-1.75)
        {\code{Analyzers.perField()}\\\scriptsize nome / contenuto};

  % --- ricerca (riga in basso)
  \node[store] (console) at (0,-3.5)   {console\\\scriptsize riga digitata};
  \node[box]   (qb)      at (3.7,-3.5) {\code{QueryBuilder}\\\scriptsize QueryParser};
  \node[box]   (fs)      at (7.5,-3.5) {\code{FileSearcher}\\\scriptsize IndexSearcher};
  \node[store] (res)     at (11,-3.5)  {risultati\\\scriptsize top-5};
  \draw[arr] (console) -- (qb);
  \draw[arr] (qb) -- node[above, font=\scriptsize]{Query} (fs);
  \draw[arr] (fs) -- (res);
  \draw[arr, dashed] (index.south) -- (index.south |- 0,-2.45) -| (fs.north)
        node[pos=0.25, above, font=\scriptsize]{lettura};

  \draw[arr, warm] (an.north) -- (iw.south west);
  \draw[arr, warm] (an.south) -- (qb.north east);

  \node[font=\small\bfseries, anchor=west] at (-1,0.95) {Indicizzazione};
  \node[font=\small\bfseries, anchor=west] at (-1,-4.45) {Ricerca};
\end{tikzpicture}
\caption{I due flussi del programma. Lo stesso analyzer per campo (in rosso) viene
usato sia per scrivere l'indice sia per analizzare la query.}
\label{fig:architettura}
\end{figure}

% ===========================================================================
\section{Dataset}
% ===========================================================================

Sono stati usati due corpora. Il \textbf{corpus piccolo} (202 file) è servito per lo
sviluppo e per le 21 query di test; il \textbf{corpus grande} (2\,810 file) è servito
per misurare i tempi su molti documenti e per valutare precision e recall. Entrambi
sono testi di Wikipedia in italiano e in inglese, scaricati con le API pubbliche di
Wikipedia; i testi sono distribuiti con licenza CC~BY-SA. Non sono un benchmark
pubblico, ma servono a provare l'indicizzazione su testo reale.

\subsection{Corpus piccolo (sviluppo e test)}
\label{sec:corpus-piccolo}

Il corpus piccolo contiene
\textbf{200 articoli di Wikipedia scelti a caso}, 100 in italiano e 100 in inglese,
scaricati tramite le API di Wikipedia. Sono stati scartati gli articoli troppo
corti, sotto i 5000 caratteri di testo. A questi si aggiungono due piccoli file
scritti a mano (\code{example1.txt}, \code{example2.txt}). I testi di Wikipedia
sono distribuiti con licenza CC~BY-SA.

Ogni file contiene il titolo dell'articolo nella prima riga e poi il testo in
chiaro. Il nome del file è \code{<lingua>\_<Titolo\_con\_underscore>.txt}
(ad esempio \code{en\_History\_of\_martial\_arts.txt}), così il campo \field{nome}
contiene parole significative da cercare.

\begin{table}[!htbp]
\centering
\begin{tabular}{@{}lrrr@{}}
\toprule
& \textbf{File} & \textbf{Parole} & \textbf{Dimensione} \\
\midrule
Italiano & 100 & 215\,394 & 1,44 MB \\
Inglese  & 100 & 209\,421 & 1,32 MB \\
Esempi   &   2 &       23 & < 1 KB \\
\midrule
\textbf{Totale} & \textbf{202} & \textbf{424\,838} & \textbf{2,76 MB} \\
\bottomrule
\end{tabular}
\caption{Composizione del corpus. La dimensione mediana di un file è di circa
10 KB, la massima di 103 KB.}
\label{tab:dataset}
\end{table}

\subsection{Corpus grande (tempi e valutazione)}
\label{sec:corpus-grande}

Il corpus grande è stato costruito con lo script \code{tools/build\_corpus.py}
(Python, solo libreria standard), incluso nel repository, in modo che chiunque possa
rigenerarlo. Il procedimento è il seguente:
\begin{enumerate}
  \item \textbf{Articoli per argomento.} Per ciascuna lingua sono state scelte 14
        categorie di Wikipedia (terremoti, battaglie della seconda guerra mondiale,
        guerre napoleoniche, squadre di calcio italiane, teoremi, pittori italiani,
        composti chimici, vulcani, fiumi italiani, imperatori romani, malattie
        infettive, opere liriche, aerei militari e linguaggi di programmazione). Dalle
        categorie, scendendo nelle sottocategorie fino a due livelli, si prendono fino
        a 80 articoli per argomento e per lingua. Così il corpus contiene gruppi di
        documenti sullo stesso tema, con parole che ricorrono in molti file e parole
        rare.
  \item \textbf{Articoli casuali.} Altri 400 articoli casuali per lingua (369 italiani
        e 346 inglesi effettivamente salvati) fanno da rumore di fondo e sono su
        qualunque argomento.
  \item \textbf{Filtri e formato.} Sono scartati gli articoli sotto i 3\,000 caratteri.
        Ogni file contiene il titolo e il testo in chiaro e si chiama
        \code{<lingua>\_<Titolo\_con\_underscore>.txt}, come nel corpus piccolo.
  \item \textbf{Manifest.} Il file \code{data/corpus-manifest.csv} elenca per ogni
        articolo lingua, titolo, nome del file, argomento di provenienza, lunghezza e
        categorie di Wikipedia. Descrive com'è fatto il corpus.
\end{enumerate}

\begin{table}[!htbp]
\centering
\begin{tabular}{@{}lrrr@{}}
\toprule
& \textbf{File} & \textbf{Parole} & \textbf{Dimensione} \\
\midrule
Italiano & 1\,401 & 2\,419\,562 & 16,2 MB \\
Inglese  & 1\,409 & 2\,868\,253 & 18,2 MB \\
\midrule
\textbf{Totale} & \textbf{2\,810} & \textbf{5\,287\,815} & \textbf{34,4 MB} \\
\bottomrule
\end{tabular}
\caption{Composizione del corpus grande. La dimensione mediana di un file è di circa
6,9 KB, la massima di 145 KB.}
\label{tab:dataset-grande}
\end{table}

\begin{table}[!htbp]
\centering
\small
\begin{tabular}{@{}lrrlrr@{}}
\toprule
\textbf{Argomento} & \textbf{it} & \textbf{en} & \textbf{Argomento} & \textbf{it} & \textbf{en} \\
\midrule
Terremoti & 74 & 80 & Fiumi & 54 & 27 \\
Battaglie 2ª g.\,m. & 80 & 80 & Imperatori romani & 80 & 80 \\
Guerre napoleoniche & 80 & 80 & Malattie infettive & 80 & 80 \\
Squadre di calcio & 80 & 80 & Opere liriche & 80 & 80 \\
Teoremi & 80 & 80 & Aerei militari & 80 & 80 \\
Pittori & 80 & 80 & Linguaggi di prog. & 80 & 80 \\
Composti chimici & 41 & 80 & Articoli casuali & 369 & 346 \\
Vulcani & 63 & 76 & & & \\
\bottomrule
\end{tabular}
\caption{Numero di articoli per argomento e lingua. Alcune categorie italiane (chimica,
fiumi, vulcani, terremoti) contengono meno di 80 articoli abbastanza lunghi.}
\label{tab:argomenti}
\end{table}

% ===========================================================================
\section{Analyzer scelti}
% ===========================================================================

I due campi hanno caratteristiche molto diverse, quindi ricevono analyzer diversi,
combinati con un \code{PerFieldAnalyzerWrapper}. L'analyzer di default del wrapper
(\code{StandardAnalyzer}) non viene mai usato, perché i campi sono solo due.
Gli analyzer sono stati scelti in due passi: una prima versione, che il confronto con
\code{grep} ha mostrato inadeguata, e la versione definitiva.

\subsection{Prima versione e problemi riscontrati}
\label{sec:prima}

La prima versione era la scelta più immediata:
\begin{itemize}
  \item \field{contenuto}: \code{StandardAnalyzer} con l'\emph{unione} delle stopword
        italiane e inglesi di Lucene;
  \item \field{nome}: \code{StandardTokenizer}, \code{LowerCaseFilter} e
        \code{WordDelimiterGraphFilter}.
\end{itemize}
Sembrava ragionevole e le prime prove funzionavano. Confrontando però il numero di
file trovati con quello di \code{grep -ilw} sui file originali, sono emersi
\textbf{molti file che il sistema non trovava}:

\begin{table}[!htbp]
\centering
\begin{tabular}{@{}lrrr@{}}
\toprule
\textbf{Termine cercato} & \textbf{\code{grep} (reale)} & \textbf{Prima versione} &
\textbf{Versione definitiva} \\
\midrule
\code{uomo}    & 19 & \textcolor{warm}{10}  & 19 \\
\code{italia}  & 38 & \textcolor{warm}{28}  & 38 \\
\code{impero}  & 10 & \textcolor{warm}{5}   & 10 \\
\code{anno}    & 55 & \textcolor{warm}{43}  & 53 \\
\code{king}    & 21 & \textcolor{warm}{18}  & 21 \\
\code{citta} (senza accento) & 34 & \textcolor{warm}{0} & 34 \\
\code{era}     & 95 & \textcolor{warm}{nessun termine utile} & 95 \\
\bottomrule
\end{tabular}
\caption{File trovati per alcuni termini, con la prima versione degli analyzer e con
quella definitiva, a confronto con il numero reale ottenuto con \code{grep}.}
\label{tab:confronto}
\end{table}

\begin{figure}[!htbp]
\centering
\begin{tikzpicture}[x=1cm, y=0.065cm]
  \draw[->] (-0.2,0) -- (13.2,0);
  \draw[->] (-0.2,0) -- (-0.2,66) node[above, font=\small] {file trovati};
  \foreach \y in {0,20,40,60}
    \draw (-0.25,\y) -- (-0.15,\y) node[left=1mm, font=\scriptsize] {\y};
  \foreach \n/\g/\p/\d [count=\i from 0] in
      {uomo/19/10/19, italia/38/28/38, impero/10/5/10, anno/55/43/53,
       king/21/18/21, citta/34/0/34} {
    \pgfmathsetmacro{\x}{\i*2.15+0.2}
    \fill[gray!55]  (\x,0)       rectangle (\x+0.55,\g);
    \fill[warm]     (\x+0.6,0)   rectangle (\x+1.15,\p);
    \fill[accent]   (\x+1.2,0)   rectangle (\x+1.75,\d);
    \node[above, font=\tiny] at (\x+0.275,\g) {\g};
    \node[above, font=\tiny] at (\x+0.875,\p) {\p};
    \node[above, font=\tiny] at (\x+1.475,\d) {\d};
    \node[below, font=\scriptsize\ttfamily] at (\x+0.875,-1) {\n};
  }
  \fill[gray!55] (9.0,58) rectangle (9.5,62);
  \node[right, font=\scriptsize] at (9.6,60) {grep (reale)};
  \fill[warm] (9.0,50) rectangle (9.5,54);
  \node[right, font=\scriptsize] at (9.6,52) {prima versione};
  \fill[accent] (9.0,42) rectangle (9.5,46);
  \node[right, font=\scriptsize] at (9.6,44) {versione definitiva};
\end{tikzpicture}
\caption{File trovati per termine: nella prima versione una parte dei file, a
volte tutti, andava persa.}
\label{fig:confronto}
\end{figure}

Le cause sono tre, tutte verificate con prove dirette sugli analyzer.

\paragraph{1. Gli apostrofi non separano le parole.}
Lo \code{StandardTokenizer} segue le regole Unicode e non spezza su un apostrofo
tra due lettere: \code{dell'uomo} resta \emph{un solo token}, e così \code{l'Italia}
e \code{nell'anno}. Nel testo italiano di Wikipedia queste forme sono continue
(\code{dell'anno} compare 28 volte, \code{l'anno} 26), quindi cercare \code{uomo}
non trovava i file in cui compare solo come \code{dell'uomo} o \code{l'uomo}:
10 file invece di 19. Lo stesso avviene per il possessivo inglese
(\code{king's} resta un token a sé), con 18 file invece di 21.

\paragraph{2. Gli accenti non sono normalizzati.}
\code{città} e \code{citta} sono termini diversi: cercando senza accento, come si
fa spesso, la prima versione trovava \textbf{0 file} contro i 34 reali.

\paragraph{3. Le stopword delle due lingue si scontrano.}
L'unione delle liste elimina da \emph{tutti} i testi le parole vuote di entrambe le
lingue. Ma 27 stopword italiane compaiono anche nel dizionario inglese del sistema, e alcune
sono parole inglesi comuni, per esempio \emph{era}, \emph{dove} e \emph{all}.
Con la prima versione la query \code{contenuto era} rispondeva ``nessun termine
utile'', pur esistendo 95 file che contengono la parola.

\paragraph{Campo \field{nome}.}
Il \code{WordDelimiterGraphFilter} separava anche lettere e numeri:
\code{example1} diventava \code{example} e \code{1}, e con l'operatore OR la query
\code{nome example1} restituiva 4 file, tre dei quali trovati solo per una parte del termine:
\code{example2} per \code{example}, \code{en\_2025\_GN1} e
\code{it\_De\_Havilland\_Canada\_DHC\_1\_Chipmunk} per il numero \code{1}.

\medskip
\noindent\textbf{Lezione appresa:} un analyzer non si valuta guardando i token di
due frasi di prova, ma confrontando il numero di risultati con una verità
indipendente (qui \code{grep}) su termini reali del corpus.

\subsection{Versione definitiva: campo \field{nome}}

\textbf{Scelta:} un \code{CustomAnalyzer} con un \code{CharFilter} che sostituisce
\code{\_} con uno spazio, poi \code{StandardTokenizer}, \code{LowerCaseFilter} e
\code{ASCIIFoldingFilter}; senza stopword e senza stemming.

\textbf{Motivazione:}
\begin{itemize}
  \item i nomi dei file usano spesso \code{\_}, \code{-} e \code{.} come
        separatori. Lo \code{StandardTokenizer} spezza su trattini e punti ma
        \emph{non} su \code{\_}, che lo standard Unicode considera parte della
        parola: il \code{CharFilter} lo trasforma in uno spazio \emph{prima} del
        tokenizer, e il problema è risolto;
  \item si è tolto il \code{WordDelimiterGraphFilter}: lettere e numeri non vengono
        più separati, quindi \code{example1} resta un unico termine;
  \item minuscolo e normalizzazione degli accenti rendono la ricerca indipendente
        da maiuscole e accenti;
  \item nei nomi anche le parole brevi (\emph{di}, \emph{the}) possono essere
        significative, e lo stemming altererebbe nomi propri e sigle, quindi non
        si usano né stopword né stemming;
  \item l'estensione \code{.txt} viene tolta prima dell'indicizzazione:
        altrimenti \code{txt} sarebbe un termine presente in ogni documento.
\end{itemize}

\subsection{Versione definitiva: campo \field{contenuto}}

\textbf{Scelta:} un analyzer composto da \code{StandardTokenizer},
\code{ElisionFilter} (articoli italiani), \code{EnglishPossessiveFilter},
\code{LowerCaseFilter} e \code{ASCIIFoldingFilter}; \textbf{senza stopword e senza
stemming}.

\textbf{Motivazione:} il corpus è bilingue.
\begin{itemize}
  \item i filtri aggiunti risolvono i primi due problemi della
        sezione~\ref{sec:prima}: l'elisione toglie l'articolo davanti all'apostrofo
        (\code{dell'uomo} $\rightarrow$ \code{uomo}), il filtro del possessivo toglie
        \code{'s} (\code{king's} $\rightarrow$ \code{king}) e la normalizzazione degli
        accenti rende uguali \code{città} e \code{citta};
  \item \textbf{nessuno stemming}: un \code{ItalianAnalyzer} o un
        \code{EnglishAnalyzer} applicherebbero lo stemming di una lingua anche ai
        testi dell'altra, storpiandone le parole;
  \item \textbf{nessuna lista di stopword}: senza riconoscere la lingua di ogni
        testo, le liste delle due lingue si scontrano (terzo problema). Il ranking
        BM25 abbassa già da solo il peso dei termini molto frequenti, quindi la perdita
        di precisione è piccola, come mostrano i test (query~12 e~18). Le
        alternative valutate, con le misure, sono nella
        sezione~\ref{sec:alternative}.
\end{itemize}

\begin{lstlisting}
// nome
CustomAnalyzer.builder()
    .addCharFilter(PatternReplaceCharFilterFactory.class,
                   "pattern", "_", "replacement", " ")
    .withTokenizer(StandardTokenizerFactory.class)
    .addTokenFilter(LowerCaseFilterFactory.class)
    .addTokenFilter(ASCIIFoldingFilterFactory.class)
    .build();

// contenuto
Tokenizer source = new StandardTokenizer();
TokenStream chain = new ElisionFilter(source, ITALIAN_ARTICLES);
chain = new EnglishPossessiveFilter(chain);
chain = new LowerCaseFilter(chain);
chain = new ASCIIFoldingFilter(chain);
\end{lstlisting}

\begin{table}[!htbp]
\centering
\begin{tabularx}{\linewidth}{@{}lXl@{}}
\toprule
\textbf{Componente} & \textbf{Cosa fa} & \textbf{Esempio} \\
\midrule
\code{PatternReplaceCharFilter} & sostituisce \code{\_} con uno spazio prima del
  tokenizer & \code{Finale\_2024} $\rightarrow$ \code{Finale 2024} \\
\code{StandardTokenizer} & divide sui confini di parola (Unicode), scarta la
  punteggiatura & \code{e-mail,} $\rightarrow$ \code{e}, \code{mail} \\
\code{ElisionFilter} & toglie l'articolo davanti all'apostrofo, per gli articoli
  italiani dell'elenco & \code{dell'uomo} $\rightarrow$ \code{uomo} \\
\code{EnglishPossessiveFilter} & toglie il possessivo inglese &
  \code{dog's} $\rightarrow$ \code{dog} \\
\code{LowerCaseFilter} & porta in minuscolo & \code{Roma} $\rightarrow$ \code{roma} \\
\code{ASCIIFoldingFilter} & toglie gli accenti & \code{città} $\rightarrow$
  \code{citta} \\
\bottomrule
\end{tabularx}
\caption{I componenti degli analyzer definitivi.}
\label{tab:filtri}
\end{table}

\subsection{Alternative valutate}
\label{sec:alternative}

Prima di eliminare le stopword dalla versione definitiva sono state confrontate
quattro varianti dello stesso analyzer di \field{contenuto}, cambiando solo l'elenco
delle parole eliminate (tabella~\ref{tab:stopword}). Per ciascuna è stato costruito
l'indice dei 202 file e ne sono stati misurati la dimensione su disco e il tempo
di indicizzazione.

\begin{table}[!htbp]
\centering
\begin{tabularx}{\linewidth}{@{}Xrrr@{}}
\toprule
\textbf{Variante} & \textbf{Parole eliminate} & \textbf{Indice} & \textbf{Tempo} \\
\midrule
Nessuna stopword (\textbf{scelta finale}) & 0 & 1299 KB & 148 ms \\
Solo articoli (\emph{il, lo, la, gli, le, un, uno, una, the, a, an}) & 11 & 1261 KB & 132 ms \\
Lista selettiva (tutte tranne 8 parole ambigue) & 302 & 1133 KB & 124 ms \\
Lista completa (italiano + inglese, prima versione) & 310 & 1128 KB & 123 ms \\
\bottomrule
\end{tabularx}
\caption{Effetto dell'elenco di stopword su dimensione dell'indice e tempo di
indicizzazione (202 file).}
\label{tab:stopword}
\end{table}

\paragraph{Lista completa.} Riduce l'indice del 13\% ma, come visto nella
sezione~\ref{sec:prima}, elimina da tutti i testi parole che nell'altra lingua sono
significative.

\paragraph{Solo articoli.} Elimina solo le parole più innocue, ma il vantaggio è
minimo (3\%): lo spazio è occupato soprattutto da preposizioni e verbi ausiliari
(\emph{di}, \emph{che}, \emph{avere}, \emph{essere}).

\paragraph{Lista selettiva.} Per escludere le parole ambigue si è misurata, per ogni
stopword, la frequenza (ogni 10\,000 parole) nei testi italiani e in quelli inglesi.
Solo 7 parole risultano frequenti nella lingua sbagliata (\emph{all}, \emph{i},
\emph{era}, \emph{ed}, \emph{c}, \emph{io}, \emph{mi}), alle quali si è aggiunta
\emph{dove}. Con queste esclusioni le prestazioni sono praticamente quelle della
lista completa (1133 contro 1128~KB). Il metodo ha però un limite: restano
eliminate parole che nei testi inglesi esistono davvero, anche se meno frequenti
che in italiano (tabella~\ref{tab:ambigue}). Ad esempio \code{non} compare in 29 file
inglesi, in gran parte perché il tokenizer spezza i composti come \emph{non-profit};
chi cerca \code{non} in un testo inglese non troverebbe il file. Inoltre l'elenco
dipende dal corpus e andrebbe rifatto per testi diversi.

\begin{table}[!htbp]
\centering
\begin{tabular}{@{}lrr@{}}
\toprule
\textbf{Stopword italiana} & \textbf{File inglesi} & \textbf{File italiani} \\
\midrule
\code{non}  & 29 &  86 \\
\code{per}  & 22 & 100 \\
\code{come} & 16 &  93 \\
\bottomrule
\end{tabular}
\caption{Stopword italiane che la lista selettiva elimina (non risultano ambigue
nella misura di frequenza) ma che compaiono comunque in molti file inglesi.}
\label{tab:ambigue}
\end{table}

\paragraph{Riconoscimento della lingua.} L'approccio corretto sarebbe applicare
l'analyzer (e la lista di stopword) della lingua di ciascun documento. La lingua si
riconosce bene anche senza librerie: contando nel testo alcune parole funzionali
(\emph{di, e, il, la, che, per, \ldots} contro \emph{the, of, and, to, in, \ldots}) e
scegliendo la lingua con la percentuale più alta, sui 200 articoli di Wikipedia la
lingua è stata indovinata in \textbf{200 casi su 200}, con un margine minimo di 6
punti percentuali tra le due lingue. Usarla in Lucene è però più complesso: l'analyzer
è associato al campo e non al documento, quindi servirebbero un campo per lingua
(\code{contenuto\_it}, \code{contenuto\_en}) oppure un \code{TokenStream} costruito
per ogni documento, e le query dovrebbero essere analizzate con entrambi gli
analyzer, perché la lingua di una query di una o due parole non è affidabile.

\paragraph{Conclusione.} Si è scelto di \textbf{non usare stopword}. Il vantaggio
della lista completa è piccolo (13\% di spazio, pochi millisecondi su 202 file, un
guadagno che diventa rilevante solo su grandi collezioni), mentre il costo è una
ricerca scorretta; la lista selettiva ridurrebbe il costo ma non lo azzera ed è
legata al corpus; il riconoscimento della lingua è la soluzione più corretta, ma
richiede molto codice in più per un risultato che con l'analyzer attuale è già
verificato (i totali coincidono con \code{grep}). Il ranking BM25 dà già un peso
molto basso ai termini molto frequenti, quindi la precisione non ne risente in modo
sensibile. Il riconoscimento della lingua resta come possibile estensione
(sezione~\ref{sec:limiti}).

\subsection{Verifica}

La figura~\ref{fig:token} riporta i token realmente prodotti dagli analyzer
definitivi (output del programma di prova).

\begin{figure}[!htbp]
\centering
\begin{tikzpicture}[node distance=2mm]
  \node[anchor=west, font=\small] (in1) at (0,0)
       {\field{nome}: \code{Relazione\_Finale-2024}};
  \node[tok, anchor=west] (a1) at (7.6,0) {relazione};
  \node[tok, right=of a1] (a2) {finale};
  \node[tok, right=of a2] (a3) {2024};
  \draw[arr] (in1.east) -- (a1.west);

  \node[anchor=west, font=\small] (in2) at (0,-1)
       {\field{nome}: \code{Città\_di\_Perché}};
  \node[tok, anchor=west] (b1) at (7.6,-1) {citta};
  \node[tok, right=of b1] (b2) {di};
  \node[tok, right=of b2] (b3) {perche};
  \draw[arr] (in2.east) -- (b1.west);

  \node[anchor=west, font=\small, align=left] (in3) at (0,-2.2)
       {\field{contenuto}: \code{L'Italia e dell'uomo}\\
        \hspace*{19mm}\code{nell'anno 1861}};
  \node[tok, anchor=west] (c1) at (7.6,-2.2) {italia};
  \node[tok, right=of c1] (c2) {e};
  \node[tok, right=of c2] (c3) {uomo};
  \node[tok, right=of c3] (c4) {anno};
  \node[tok, right=of c4] (c5) {1861};
  \draw[arr] (in3.east) -- (c1.west);

  \node[anchor=west, font=\small] (in4) at (0,-3.4)
       {\field{contenuto}: \code{The dog's bone, perché?}};
  \node[tok, anchor=west] (d1) at (7.6,-3.4) {the};
  \node[tok, right=of d1] (d2) {dog};
  \node[tok, right=of d2] (d3) {bone};
  \node[tok, right=of d3] (d4) {perche};
  \draw[arr] (in4.east) -- (d1.west);
\end{tikzpicture}
\caption{Token prodotti dagli analyzer definitivi. Il carattere \code{\_} separa le
parole, gli articoli elisi e il possessivo vengono tolti, accenti e maiuscole
normalizzati; le parole comuni (\emph{e}, \emph{the}) restano.}
\label{fig:token}
\end{figure}

\subsection{Limiti noti}
\begin{itemize}
  \item \textbf{Nessuno stemming:} cercare \emph{libro} non trova \emph{libri}.
        È il prezzo della scelta multilingua.
  \item \textbf{Nessuna stopword:} una query fatta solo di parole comuni, come
        \code{contenuto il the}, trova quasi tutti i file (199 su 202, e 2\,803 su 2\,810 nel
        corpus grande); il ranking
        li mette comunque in ordine di rilevanza, ma l'utente non viene avvisato.
  \item \textbf{Elenco degli articoli incompleto:} \code{quell'anno} (7 occorrenze)
        non viene trattato perché \code{quell} non è nell'elenco, e per questo
        \code{anno} trova 53 file invece di 55 (nel corpus grande la valutazione trova
        casi analoghi: \emph{quest'uomo}, \emph{pover'uomo}, \emph{quand'era},
        \emph{com'era}).
  \item \textbf{Indirizzi web nel testo:} lo \code{StandardTokenizer} non spezza una
        parola su un punto tra due lettere, per cui in un testo con
        \code{volcano.und} il termine \emph{volcano} non viene trovato.
  \item \textbf{Numeri nei nomi:} \code{example1} è un unico termine, quindi la
        query \code{nome example} non trova più il file \code{example1}. È il
        rovescio della scelta fatta per evitare i risultati spuri.
\end{itemize}

% ===========================================================================
\section{Indicizzazione}
% ===========================================================================

Si crea un \code{Document} Lucene per ogni file, con due campi
(tabella~\ref{tab:campi}). La cartella viene visitata ricorsivamente e vengono presi
i file con estensione \code{.txt}, senza distinguere maiuscole e minuscole.
L'indice è ricreato da zero a ogni esecuzione (\code{OpenMode.CREATE}), così
rieseguire il programma non duplica i documenti. Un file illeggibile (ad esempio
non codificato in UTF-8) viene saltato con un avviso e contato a parte, senza
interrompere il lavoro. Alla fine viene fatto il \code{commit()} e vengono stampati
il numero di file e il tempo impiegato.

\begin{table}[!htbp]
\centering
\begin{tabularx}{\linewidth}{@{}llcX@{}}
\toprule
\textbf{Campo} & \textbf{Tipo} & \textbf{Stored} & \textbf{Perché} \\
\midrule
\field{nome} & \code{TextField} & sì & serve per stampare il nome nei risultati \\
\field{contenuto} & \code{TextField} & no & serve solo per cercare: non salvarlo
  riduce lo spazio occupato dall'indice \\
\bottomrule
\end{tabularx}
\caption{I campi del documento Lucene.}
\label{tab:campi}
\end{table}

\subsection{Risultati sul corpus piccolo}

Sul corpus piccolo, su cinque esecuzioni consecutive con gli analyzer definitivi sono stati indicizzati
\textbf{202 file}, con \textbf{0 file saltati}, in circa \textbf{0,6 secondi}
(figura~\ref{fig:tempi}). L'indice occupa circa 1,3 MB, contro i 2,8 MB di testo.
La prima versione degli analyzer impiegava circa 0,5 secondi (tre esecuzioni:
565, 466 e 477 ms) e produceva un indice di circa 1,2 MB: i filtri aggiunti e
l'assenza di stopword costano poco, dello stesso ordine di grandezza.

\begin{figure}[!htbp]
\centering
\begin{tikzpicture}[x=1.3cm, y=0.008cm]
  % assi
  \draw[->] (0.3,0) -- (5.9,0);
  \draw[->] (0.3,0) -- (0.3,800) node[above, font=\small] {ms};
  \foreach \y in {0,200,400,600}
    \draw (0.27,\y) -- (0.33,\y) node[left=1mm, font=\scriptsize] {\y};
  % barre
  \foreach \x/\v in {1/696, 2/691, 3/577, 4/489, 5/502} {
    \fill[accent] (\x-0.3,0) rectangle (\x+0.3,\v);
    \node[below, font=\small, text=white] at (\x,\v) {\v};
    \node[below, font=\small] at (\x,0) {run \x};
  }
  % media
  \draw[warm, thick, dashed] (0.4,591) -- (5.5,591)
       node[right, font=\scriptsize, text=warm] {media 591};
\end{tikzpicture}
\caption{Tempo totale di indicizzazione dei 202 file in cinque esecuzioni.}
\label{fig:tempi}
\end{figure}

\noindent\textbf{Ambiente di misura:} MacBook Pro (2020) con chip
Apple M1 a 8 core (4 prestazioni e 4 efficienza), 8~GB di RAM, macOS~26.6,
OpenJDK~21.0.10.


\noindent Il tempo misurato comprende la lettura dei file e l'analisi del testo
ed è preso in una JVM appena avviata, senza riscaldamento: le prime esecuzioni
sono le più lente. Per questo conviene considerare l'ordine di grandezza
(circa 0,5--0,7 s) più del valore singolo. Il costo per file (2--3 ms) è però
sovrastimato, perché su pochi file pesa molto il tempo fisso di avvio di JVM e Lucene:
lo mostra la misura sul corpus grande.

\subsection{Tempi al crescere del corpus}
\label{sec:tempi-grande}

Per capire come crescono i tempi, lo script \code{tools/benchmark\_indexing.py} costruisce
sottoinsiemi casuali (con seme fisso, quindi sempre uguali) del corpus grande da 200,
500, 1\,000, 2\,000 e 2\,810 file e lancia l'indicizzatore cinque volte per ciascuno,
ogni volta in una nuova JVM. Nella tabella~\ref{tab:scala} si riporta la
\textbf{mediana}: una delle cinque esecuzioni sul corpus completo è durata 3\,289~ms
contro circa 2\,150~ms delle altre (probabilmente un disturbo del sistema) e la
mediana non ne risente.

\begin{table}[!htbp]
\centering
\begin{tabular}{@{}rrrrr@{}}
\toprule
\textbf{File} & \textbf{Testo (MB)} & \textbf{Indice (MB)} & \textbf{Mediana (ms)} & \textbf{ms/file} \\
\midrule
200   & 2,1  & 1,05  &  461 & 2,31 \\
500   & 5,7  & 2,52  &  705 & 1,41 \\
1\,000 & 12,0 & 4,86  & 1\,013 & 1,01 \\
2\,000 & 23,9 & 9,63  & 1\,664 & 0,83 \\
2\,810 & 34,4 & 13,73 & 2\,326 & 0,83 \\
\bottomrule
\end{tabular}
\caption{Tempo di indicizzazione al crescere del corpus (mediana di cinque esecuzioni).}
\label{tab:scala}
\end{table}

\begin{figure}[!htbp]
\centering
\begin{tikzpicture}[x=0.0045cm, y=0.0022cm]
  \draw[->] (0,0) -- (3050,0) node[right, font=\small] {file};
  \draw[->] (0,0) -- (0,2600) node[above, font=\small] {ms};
  \foreach \x in {500,1000,1500,2000,2500,3000}
    \draw (\x,0) -- (\x,-40) node[below, font=\scriptsize] {\x};
  \foreach \y in {500,1000,1500,2000,2500}
    \draw (0,\y) -- (-30,\y) node[left, font=\scriptsize] {\y};
  \draw[accent, thick] (200,461) -- (500,705) -- (1000,1013) -- (2000,1664) -- (2810,2326);
  \foreach \x/\y in {200/461,500/705,1000/1013,2000/1664,2810/2326}
    \fill[warm] (\x,\y) circle (2.5pt);
  \node[font=\scriptsize, anchor=south east] at (2810,2326) {2\,326 ms};
\end{tikzpicture}
\caption{Tempo di indicizzazione in funzione del numero di file (mediana).}
\label{fig:scala}
\end{figure}

\noindent La crescita è quasi lineare. Con pochi file il tempo è dominato da un costo
fisso di circa 400~ms (avvio di JVM e Lucene); oltre i 2\,000 file il costo si
stabilizza a circa \textbf{0,83~ms per file}, cioè circa 15~MB di testo al secondo.
L'indice occupa circa il 40\% del testo, perché il contenuto non è memorizzato.
Indicizzare tutti i 2\,810 file richiede quindi circa 2,3 secondi.

% ===========================================================================
\section{Ricerca}
% ===========================================================================

\subsection{Sintassi della query}

Ogni riga digitata ha la forma \code{<campo> <termini>}, dove \code{<campo>} è
\field{nome} o \field{contenuto}. Il
testo tra virgolette è una frase. Ad esempio:
\begin{center}
\code{nome knitting}\\
\code{contenuto guerra mondiale}\\
\code{contenuto \dq{}world war\dq}
\end{center}

\subsection{Costruzione della query}

La parola chiave iniziale viene separata a mano (\code{split(\dq{} \dq{}, 2)}), mentre il
resto della riga è affidato al \code{QueryParser} di Lucene, configurato con il
campo scelto come \emph{campo di default} e con lo stesso analyzer usato in
indicizzazione. In questo modo le frasi tra virgolette e l'analisi dei termini
sono gestite da Lucene, con meno codice e meno casi limite rispetto a un parser
scritto a mano.

\textbf{Più termini senza virgolette sono in OR.} È il comportamento predefinito
del \code{QueryParser}: i termini senza
modificatori sono \emph{opzionali} (\code{BooleanClause.\allowbreak Occur.\allowbreak SHOULD}). I documenti
che contengono più termini ottengono comunque un punteggio maggiore. Il ranking è
quello predefinito di Lucene (BM25).

\begin{table}[!htbp]
\centering
\begin{tabular}{@{}ll@{}}
\toprule
\textbf{Riga digitata} & \textbf{Query Lucene generata} \\
\midrule
\code{nome roma} & \code{nome:roma} \\
\code{contenuto data scientist} & \code{contenuto:data contenuto:scientist} \\
\code{contenuto \dq{}data scientist\dq} & \code{contenuto:\dq{}data scientist\dq} \\
\code{contenuto Data} & \code{contenuto:data} \\
\code{nome Relazione\_Finale} & \code{nome:relazione nome:finale} \\
\bottomrule
\end{tabular}
\caption{Traduzione delle righe digitate (output reale di \code{query.toString()}).}
\label{tab:parsing}
\end{table}

\subsection{Esecuzione e risultati}

Il \code{FileSearcher} apre l'indice con \code{DirectoryReader} e
\code{IndexSearcher} e viene chiuso automaticamente a fine sessione
(\code{try-with-resources}). Per ogni query stampa il numero totale di file
trovati e i primi \textbf{5} (posizione, nome e punteggio). Il nome viene letto
dal campo memorizzato; il contenuto non si può mostrare perché non è memorizzato.

Gli errori di input sono gestutu: campo sconosciuto, testo
mancante e sintassi non valida (ad esempio virgolette non chiuse) producono un
messaggio e il ciclo prosegue. Se tutti i termini vengono eliminati
dall'analyzer (ad esempio una query fatta solo di simboli, come \code{c++}) la
query è vuota e il programma lo segnala invece di restituire zero risultati senza
spiegazione.

% ===========================================================================
\section{Test del sistema}
\label{sec:test}
% ===========================================================================

Le 21 query di test sono eseguite sul \textbf{corpus piccolo} (202 file). Sono
raccolte nel file \code{queries.txt} ed eseguite in modo
automatico dallo script \code{run-queries.sh}, che salva l'output completo in
\code{results/query-results.md}. I test sono quindi riproducibili. Per le query sul
contenuto con un termine o una frase (query 5--11 e 14--18) il numero di risultati è
stato confrontato con quello di \code{grep -il} sui file originali: in tutti i casi
coincide. Le query 14--18 sono state aggiunte per verificare i problemi della
sezione~\ref{sec:prima}.

\textbf{Perché il corpus piccolo.} Questi test sono \emph{funzionali}: mostrano il
comportamento del programma (output, top-5, frasi, casi di errore) e non dipendono
dalla dimensione del corpus. Si usa il corpus piccolo perché con 202 file il numero
di risultati si controlla facilmente a mano con \code{grep}, e perché i problemi
della sezione~\ref{sec:prima} sono stati individuati e corretti su quei file, quindi i
numeri si possono confrontare direttamente. La correttezza su larga scala, con 2\,810
file e un riferimento indipendente da Lucene, è valutata nella
sezione~\ref{sec:valutazione}.

\begin{table}[!htbp]
\centering
\small
\begin{tabularx}{\linewidth}{@{}rlXrl@{}}
\toprule
\textbf{\#} & \textbf{Query} & \textbf{Cosa verifica} & \textbf{Ris.} &
\textbf{Primo risultato} \\
\midrule
1  & \code{nome knitting} & un termine nel nome & 1 & \code{en\_Knitting} \\
2  & \code{nome history martial} & più termini nel nome (OR) & 1 & \code{en\_History\_of\_martial\ldots} \\
3  & \code{nome \dq{}martial arts\dq} & frase nel nome & 1 & \code{en\_History\_of\_martial\ldots} \\
4  & \code{nome list} & totale oltre il top-5 & 6 & \code{en\_List\_of\_formulae\ldots} \\
5  & \code{contenuto napoleone} & termine italiano & 5 & \code{it\_Castello\_di\_Saint\ldots} \\
6  & \code{contenuto earthquake} & termine inglese & 1 & \code{en\_2020\_Petrinja\ldots} \\
7  & \code{contenuto \dq{}world war\dq} & frase nel contenuto & 16 & \code{en\_Protectorate\_of\ldots} \\
8  & \code{contenuto guerra mondiale} & più termini (OR) & 44 & \code{it\_Polemiko\_Nautiko} \\
9  & \code{contenuto NAPOLEONE} & maiuscole ignorate & 5 & come la query 5 \\
10 & \code{contenuto \dq{}data scientist\dq} & frase in un solo file & 1 & \code{example1} \\
11 & \code{contenuto xyzqwerty} & nessun risultato & 0 & -- \\
12 & \code{contenuto il the} & solo parole comuni & 199 & \code{it\_Emily\_Blunt} \\
13 & \code{nome example1} & numero nel nome & 1 & \code{example1} \\
14 & \code{contenuto uomo} & elisione (\emph{dell'uomo}) & 19 & \code{it\_Belle\_Gunness} \\
15 & \code{contenuto italia} & elisione (\emph{l'Italia}) & 38 & \code{it\_Giuliano\_Ferrara} \\
16 & \code{contenuto citta} & accenti (\emph{città}) & 34 & \code{it\_Citta\_di\_Ebla} \\
17 & \code{contenuto king} & possessivo (\emph{king's}) & 21 & \code{it\_Non\_siamo\_piu\ldots} \\
18 & \code{contenuto era} & parola italiana e inglese & 95 & \code{it\_Detlev\_Peukert} \\
19 & \code{titolo knitting} & campo sconosciuto & -- & errore gestito \\
20 & \code{contenuto} & testo mancante & -- & errore gestito \\
21 & \code{contenuto \dq{}frase aperta} & virgolette non chiuse & -- & errore gestito \\
\bottomrule
\end{tabularx}
\caption{Le 21 query di test (risultati reali, con gli analyzer definitivi).}
\label{tab:query}
\end{table}

\begin{figure}[!htbp]
\centering
\begin{tikzpicture}[x=0.1cm, y=0.5cm]
  \foreach \i/\q/\v/\c in {
      1/{nome knitting}/1/accent,
      2/{nome history martial}/1/accent,
      3/{nome \dq{}martial arts\dq}/1/warm,
      4/{nome list}/6/accent,
      5/{nome example1}/1/accent,
      6/{contenuto napoleone}/5/accent,
      7/{contenuto earthquake}/1/accent,
      8/{contenuto guerra mondiale}/44/accent,
      9/{contenuto \dq{}world war\dq}/16/warm,
      10/{contenuto \dq{}data scientist\dq}/1/warm,
      11/{contenuto xyzqwerty}/0/accent,
      12/{contenuto uomo}/19/accent,
      13/{contenuto italia}/38/accent,
      14/{contenuto citta}/34/accent,
      15/{contenuto king}/21/accent,
      16/{contenuto era}/95/accent} {
    \pgfmathsetmacro{\yy}{-\i}
    \node[anchor=east, font=\scriptsize\ttfamily] at (-3,\yy) {\q};
    \fill[\c] (0,\yy-0.3) rectangle (\v,\yy+0.3);
    \node[anchor=west, font=\scriptsize] at (\v+2,\yy) {\v};
  }
  \draw (0,-0.4) -- (0,-16.6);
  \fill[accent] (70,-2.2) rectangle (75,-1.8);
  \node[anchor=west, font=\scriptsize] at (77,-2) {termini (OR)};
  \fill[warm] (70,-3.7) rectangle (75,-3.3);
  \node[anchor=west, font=\scriptsize] at (77,-3.5) {frase};
\end{tikzpicture}
\caption{Numero di file trovati dalle query di test (in rosso le frasi tra
virgolette). Sono escluse la query 9, che duplica la 5, la query 12 (199 file) e i
casi di errore 19--21.}
\label{fig:risultati}
\end{figure}

\subsection{Osservazioni}

\begin{itemize}
  \item \textbf{Problemi risolti.} Le query 14--18 danno esattamente il numero di
        file di \code{grep} (19, 38, 34, 21, 95): elisioni, accenti, possessivi e
        parole come \emph{era} non causano più file mancanti. La query~13 trova solo
        \code{example1}, senza i risultati spuri della prima versione.
  \item \textbf{OR contro frase.} \code{contenuto guerra mondiale} trova 44 file,
        cioè tutti quelli che contengono almeno una delle due parole, mentre la
        frase \code{\dq{}world war\dq} ne trova 16: la frase è molto più precisa.
        Il ranking BM25 porta comunque in cima i file che contengono entrambe le
        parole più volte.
  \item \textbf{Analyzer coerente tra indice e query.} Le query 5 e 9 danno
        esattamente gli stessi risultati e gli stessi punteggi: il minuscolo è
        applicato sia ai documenti sia alla query.
  \item \textbf{Parole comuni.} Senza stopword, la query~12 trova 199 file su 202:
        è il costo della scelta, mitigato dal fatto che BM25 dà poco peso ai termini
        molto frequenti.
  \item \textbf{Casi di errore.} Le query 19--21 non fanno terminare il
        programma: ognuna produce un messaggio comprensibile.
\end{itemize}

% ===========================================================================
\section{Valutazione con precision e recall}
\label{sec:valutazione}
% ===========================================================================

Contare i file trovati non dice se ne mancano. Per valutare la \emph{recall} serve un
riferimento indipendente da Lucene che dica quali file \emph{dovrebbero} essere
trovati. La valutazione è fatta sul corpus grande (2\,810 file) dalla classe
\code{Evaluator}.

\subsection{Metodo}

\textbf{Verità lessicale.} Per ogni query, un file è \emph{rilevante} se contiene
davvero la parola cercata (o, per una frase tra virgolette, le parole una dopo
l'altra). L'insieme dei file rilevanti è calcolato senza usare Lucene: il programma
legge ogni file, toglie gli accenti, porta il testo in minuscolo, lo divide su ogni
carattere che non sia una lettera o una cifra e controlla se i termini della query
compaiono, come farebbe \code{grep -iw}. Più termini senza virgolette sono in OR, come
nel \code{QueryParser}. Per il campo \field{nome} si fa lo stesso sul nome del file.

\textbf{Metriche.} Per ogni query, su \emph{tutti} i risultati:
\begin{gather*}
\text{precision}=\frac{\text{file trovati e rilevanti}}{\text{file trovati}},\qquad
\text{recall}=\frac{\text{file trovati e rilevanti}}{\text{file rilevanti}},\\[2pt]
F_1=\frac{2\cdot\text{precision}\cdot\text{recall}}{\text{precision}+\text{recall}}.
\end{gather*}
Se non esistono file rilevanti e non ne viene trovato nessuno la risposta è corretta
(precision e recall pari a 1); se ne esistono ma non ne viene trovato nessuno la
precision è posta a 0. Si riporta la media sulle query (media macro).

\textbf{Cosa si confronta.} Le stesse 26 query (20 sul contenuto e 6 sul nome),
con gli stessi file, sono eseguite con due configurazioni: la \emph{prima versione}
degli analyzer (sezione~\ref{sec:prima}) e la \emph{versione definitiva}. Le query
comprendono termini italiani e inglesi, frasi, parole con elisione o accento, parole
comuni e casi sul nome.

\textbf{Controllo dell'oracolo.} Poiché la verità è scritta da noi, i suoi insiemi
sono stati ricalcolati con altri due strumenti (\code{perl} per minuscolo e accenti,
\code{grep -E} per la ricerca) con lo script \code{tools/crosscheck\_oracle.py}: i
file rilevanti coincidono \textbf{su tutte le 26 query, file per file}.

\subsection{Risultati}

\begin{table}[!htbp]
\centering
\footnotesize
\begin{tabularx}{\linewidth}{@{}Xrrrrrrr@{}}
\toprule
& & \multicolumn{3}{c}{\textbf{Prima versione}} & \multicolumn{3}{c}{\textbf{Definitiva}} \\
\cmidrule(lr){3-5}\cmidrule(lr){6-8}
\textbf{Query} & \textbf{Ril.} & \textbf{Trov.} & \textbf{P} & \textbf{R} & \textbf{Trov.} & \textbf{P} & \textbf{R} \\
\midrule
\code{contenuto napoleone} & 83 & 83 & 1,00 & 1,00 & 83 & 1,00 & 1,00 \\
\code{contenuto napoleon} & 100 & 83 & 1,00 & 0,83 & 100 & 1,00 & 1,00 \\
\code{contenuto terremoto} & 109 & 109 & 1,00 & 1,00 & 109 & 1,00 & 1,00 \\
\code{contenuto earthquake} & 140 & 140 & 1,00 & 1,00 & 140 & 1,00 & 1,00 \\
\code{contenuto vulcano} & 74 & 74 & 1,00 & 1,00 & 74 & 1,00 & 1,00 \\
\code{contenuto volcano} & 100 & 99 & 1,00 & 0,99 & 99 & 1,00 & 0,99 \\
\code{contenuto \dq{}world war\dq{}} & 349 & 349 & 1,00 & 1,00 & 349 & 1,00 & 1,00 \\
\code{contenuto \dq{}seconda guerra mondiale\dq{}} & 186 & 186 & 1,00 & 1,00 & 186 & 1,00 & 1,00 \\
\code{contenuto \dq{}roman empire\dq{}} & 119 & 119 & 1,00 & 1,00 & 119 & 1,00 & 1,00 \\
\code{contenuto \dq{}machine learning\dq{}} & 12 & 12 & 1,00 & 1,00 & 12 & 1,00 & 1,00 \\
\code{contenuto impero} & 209 & 84 & 1,00 & 0,40 & 209 & 1,00 & 1,00 \\
\code{contenuto uomo} & 211 & 133 & 1,00 & 0,63 & 209 & 1,00 & 0,99 \\
\code{contenuto italia} & 518 & 427 & 1,00 & 0,82 & 518 & 1,00 & 1,00 \\
\code{contenuto citta} & 535 & 2 & 1,00 & 0,00 & 535 & 1,00 & 1,00 \\
\code{contenuto king} & 273 & 257 & 1,00 & 0,94 & 273 & 1,00 & 1,00 \\
\code{contenuto era} & 1030 & 0 & 0,00 & 0,00 & 1027 & 1,00 & 1,00 \\
\code{contenuto anno} & 540 & 382 & 1,00 & 0,71 & 522 & 1,00 & 0,97 \\
\code{contenuto NAPOLEONE} & 83 & 83 & 1,00 & 1,00 & 83 & 1,00 & 1,00 \\
\code{contenuto teorema theorem} & 152 & 152 & 1,00 & 1,00 & 152 & 1,00 & 1,00 \\
\code{contenuto il the} & 2803 & 0 & 0,00 & 0,00 & 2803 & 1,00 & 1,00 \\
\code{nome napoleone} & 2 & 2 & 1,00 & 1,00 & 2 & 1,00 & 1,00 \\
\code{nome roma} & 7 & 7 & 1,00 & 1,00 & 7 & 1,00 & 1,00 \\
\code{nome battle} & 15 & 15 & 1,00 & 1,00 & 15 & 1,00 & 1,00 \\
\code{nome \dq{}world war\dq{}} & 10 & 10 & 1,00 & 1,00 & 10 & 1,00 & 1,00 \\
\code{nome example1} & 0 & 10 & 0,00 & 1,00 & 0 & 1,00 & 1,00 \\
\code{nome vesuvio} & 1 & 1 & 1,00 & 1,00 & 1 & 1,00 & 1,00 \\
\midrule
\textbf{Media} & & & 0,88 & 0,82 & & 1,00 & 1,00 \\
\bottomrule
\end{tabularx}
\caption{Precision (P) e recall (R) sul corpus grande, per ogni query. \emph{Ril.} è il
numero di file rilevanti per la verità lessicale, \emph{Trov.} il numero di file
trovati da Lucene. Nell'ultima riga la media macro sulle 26 query.}
\label{tab:valutazione}
\end{table}

La media macro passa, dalla prima versione alla definitiva, da precision 0,88 e recall
0,82 (F$_1$ 0,81) a precision 1,00 e recall 1,00 (F$_1$ 1,00), con arrotondamento a due cifre.
I casi più chiari sono:
\begin{itemize}
  \item \code{contenuto citta}: la prima versione trova solo 2 dei 535 file
        (recall 0,00) perché non normalizza gli accenti, la definitiva li trova tutti;
  \item \code{contenuto era}: 0 su 1\,030, perché \emph{era} è una stopword italiana;
        la definitiva trova 1\,027 file su 1\,030;
  \item \code{contenuto impero}, \code{uomo}, \code{italia}: la prima versione perde
        i file in cui la parola segue un articolo con apostrofo (\emph{dell'impero},
        \emph{l'uomo}, \emph{l'Italia}), con recall rispettivamente 0,40, 0,63 e
        0,82; \code{napoleon} (recall 0,83) perde i file che contengono solo il
        possessivo inglese \emph{Napoleon's};
  \item \code{nome example1}: la prima versione restituisce 10 file non pertinenti,
        perché \code{WordDelimiterGraphFilter} separa \emph{example1} in
        \emph{example} e \emph{1} e cerca il numero 1 in molti nomi; nel corpus non
        esiste nessun file con quel nome, quindi la risposta corretta è zero
        risultati, che la versione definitiva dà.
\end{itemize}
La precision della versione definitiva è 1,00 su tutte le query: ogni file trovato
contiene davvero la parola.

\subsection{Limiti della valutazione}

La versione definitiva non è perfetta: perde da 1 a 18 file su alcune query. Aprendo
i file corrispondenti si trovano due cause. La prima sono le \textbf{elisioni non
presenti nell'elenco degli articoli}: \emph{quest'uomo}, \emph{pover'uomo},
\emph{quand'era}, \emph{com'era}, \emph{quell'anno}; così \code{uomo} trova 209 file
su 211, \code{era} 1\,027 su 1\,030 e \code{anno} 522 su 540. La seconda sono gli
\textbf{indirizzi web} come \code{volcano.und}, che lo \code{StandardTokenizer} non
spezza. Entrambe sono descritte tra i limiti.

\textbf{Che cosa misura questa valutazione.} La verità lessicale controlla che
l'indicizzazione sia corretta, cioè che non vengano persi file che contengono
davvero la parola; non misura se i risultati siano \emph{utili} per chi cerca (un
articolo che nomina Napoleone una volta sola è comunque rilevante). Inoltre la verità
usa regole simili a quelle dell'analyzer (minuscolo, niente accenti, separazione sui
simboli), quindi l'accordo quasi totale dimostra che l'analyzer applica
correttamente queste regole e non che le regole siano le uniche possibili. Il
confronto tra le due versioni resta valido perché entrambe sono misurate con lo
stesso riferimento. Per questo motivo non si riporta la precision@5: con questa
verità ogni file trovato è rilevante, e valutare l'ordine dei risultati richiederebbe
una verità tematica (per esempio le categorie di Wikipedia, già presenti nel
manifest), non realizzata.

% ===========================================================================
\section{Limiti e possibili miglioramenti}
\label{sec:limiti}
% ===========================================================================

\begin{itemize}
  \item \textbf{Stemming e stopword per lingua:} riconoscere la lingua di ogni file
        e usare l'analyzer della lingua giusta migliorerebbe il richiamo
        (\emph{libro}/\emph{libri}) e permetterebbe di reintrodurre le stopword
        senza gli scontri tra le due lingue. La rilevazione è già fattibile
        (sezione~\ref{sec:alternative}); servirebbero un campo per lingua e
        l'analisi della query con entrambi gli analyzer.
  \item \textbf{Elisioni:} completare l'elenco degli articoli (\emph{quell'},
        \emph{quest'}, \emph{quand'}, \emph{com'}, \emph{pover'}, \emph{sant'},
        \ldots) per recuperare i file mancanti nelle query su \code{anno},
        \code{uomo} ed \code{era}, che la valutazione sul corpus grande ha
        individuato.
  \item \textbf{Indirizzi web nel testo:} un tokenizer che spezzi anche sul punto
        tra lettere (o un filtro dedicato) recupererebbe casi come \code{volcano.und}.
  \item \textbf{Verità tematica:} per valutare anche l'ordine dei risultati
        (precision@5, MAP) si potrebbero usare le categorie di Wikipedia già
        presenti nel manifest come insieme dei documenti rilevanti per un argomento.
  \item \textbf{Numeri nei nomi:} \code{nome example} non trova \code{example1}.
        Un filtro che separi lettere e numeri \emph{e} conservi il termine
        originale (opzione \code{preserveOriginal} del
        \code{WordDelimiterGraphFilter}) coprirebbe entrambi i casi.
  \item \textbf{AND come operatore predefinito:} renderebbe più precise le query
        con più termini, ma rischia di non trovare nulla; andrebbe valutato su un
        insieme di query con risultati attesi (ad esempio con la precision@5).
  \item \textbf{Caratteri speciali:} il \code{QueryParser} interpreta ancora la
        sintassi di Lucene (\code{campo:valore}, \code{+}, \code{-}, \code{*}),
        per cui \code{nome foo:bar} cerca nel campo \code{foo}. Un parser scritto a
        mano o l'escape dei caratteri speciali limiterebbero la sintassi a quella
        richiesta.
  \item \textbf{Test automatici:} i casi della tabella~\ref{tab:query} potrebbero
        diventare test JUnit.
\end{itemize}

% ===========================================================================
\section{Come riprodurre i risultati}
% ===========================================================================

\begin{lstlisting}[language=bash]
./index.sh                 # indicizza data/docs nella cartella index/
./search.sh                # ricerca interattiva (exit per uscire)
./run-queries.sh           # esegue queries.txt e salva results/query-results.md
\end{lstlisting}

\noindent Per rigenerare il corpus grande, misurare i tempi e rifare la valutazione:
\begin{lstlisting}[language=bash]
python3 tools/build_corpus.py          # scarica il corpus in data/corpus (ripartibile)
python3 tools/benchmark_indexing.py    # tempi su 200/500/1000/2000/tutti i file
# valutazione (Evaluator si lancia con java e il classpath di Maven):
./mvnw -q dependency:build-classpath -Dmdep.outputFile=cp.txt
java -cp "target/classes:$(cat cp.txt)" it.homework.lucene.Evaluator \
     data/corpus results/evaluation.md
python3 tools/crosscheck_oracle.py     # controllo dell'oracolo con perl e grep
\end{lstlisting}

\end{document}
